* Objetivo 
    - Conceitualizar, implementar e avaliar artefatos tecnológicos

* Conceitualizar
    - Processo de formulação a partir de um grounded truth (premissa).
    - Aqui se define as metricas e metodologias de avaliação dos artefatos

* implementar
    - Descrição e arguição das decisões de projeto isoladas em cenários que contribuem incrementalmente para a solução.
    - Apresentação dos processos, modelos e artefatos.

* Avaliar artefatos tecnológicos
    -  Apresentar resultados isolados para cada cenário e a construção da solução (instanciando os processos / arquiteturas ) e resultados para a solução implementada face as metricas e e metodologias de avaliações definidas na etapa de contextualização.





Nosso caso: 

Proposta de Grounded Truth (GT): 

Aqui deve vir da literatura e levar em consideração os elementos abaixo que foram retirado de livro - Artificial Intelligence - A Modern Approach

* An agent is anything that can be viewed as perceiving its environment through sensors and
acting upon that environment through actuators

* percept to refer to the agent’s perceptual inputs at any given instant. An
agent’s percept sequence is the complete history of everything the agent has ever perceived.

* agent’s behavior is described by the agent function that maps any given percept sequence to an action.

* Experience of the agent happens as it is plunked down in an environment, it generates a sequence of actions according to the percepts it receives. This sequence of actions causes the environment to go through a sequence of states. If the sequence is desirable, then the agent has performed well. This notion of desirability is captured by a performance measure that evaluates any given sequence of environment states.

* A task defines Performance, Environment, Actuators, Sensors.

Portano o problema que um agente resolve pode ser definido em termos de:

% incomplete here
* States: 
* Initial state: 
* Actions: 
* Transition model: 
* Goal test: 
%until here


Proposta de Avaliação:


                        Resultados de pesquisa x Atividades da Pesquisa
                        
                               Construção                    |           Avaliação
                        Objetivo      |    Artefatos         |    Métricas    |   Metodologia
Conceitos    |Adaptação dos conceitos | Revisão de literatura|Completeza e    |  Argumentação
             |que norteriam o problema| definição dos papeis:|Inteligibilidade|     Lógica
             |em ia para redes (GT)   | tarefa agente e exp  |                |
             |                        |                      |                |
Modelos      | Descrever como usar    | Metricas de avaliação| Completeza e   | Argumentação
             | as classes do unsup    | da representatividade| fidelidade com |     Lógica
             | para fazer o profiling | das classes e treino | mundo real     |
             |                        | online               |                |
             |                        |                      |                |
Método       | Demostrar a utilização | Combinatorial Data   | Consistencia   | Avaliação de  
             |com os dados que contem| Fusion model/process  |  interna       |Especialistas(AE)
             | o DDOS e usando  Combi-|                      |                |
             |  natorial DataFusion   |                      |                |
             |                        |                      |                |
Implementação|                        |                      |                | Comparar rotulação
             |                        |    ????              |                | gerada com rótulos de           
             |                        |                      |                | dataset públicos, e
             |                        |                      |                | comparar performance
             |                        |                      |  consistencia  | com baselines pre-existentes
             |                        |                      |     interna    | 


Metricas para custer

Silhouette Score	Cluster quality 
Davies-Bouldin		Similaridade entre clusters
Calinski-Harabasz	Variancia  intra clusters
Dunn Index	    	Min dist inter 
Cauchy-Schwarz         ??
Density-Based Clustering Validation ??